\documentclass[a4paper,11pt]{letter}
\usepackage[utf8]{inputenc}
\usepackage[T1]{fontenc}
\usepackage{geometry}
\geometry{left=2.5cm,right=2.5cm,top=2cm,bottom=2cm}
\usepackage{hyperref}
\hypersetup{colorlinks=true, urlcolor=blue}

\begin{document}

\begin{flushleft}
    \textbf{Ismail AISSAOUI} \\
    State Engineer in Artificial Intelligence (ESI-SBA) \\
    Email: ismail.aissaoui.pro@gmail.com \\
    Phone: +213 660 70 77 96 \\
    LinkedIn: linkedin.com/in/aissaoui-ismail-6a77a92a7 \\
    Portfolio: ismail-aissaoui.vercel.app
\end{flushleft}

\bigskip

\begin{flushright}
    To the Attention of Dr. Roufaida Laidi \\
    Section for Medical ICT, The Intervention Centre \\
    Oslo University Hospital (OUS) \\
    Oslo, Norway
\end{flushright}

\bigskip

\begin{center}
    \textbf{Subject: Application for PhD Fellowship in AI-Native Networking (SYNAPSE Project)}
\end{center}

\bigskip

Dear Dr. Laidi and the Selection Committee,

As a State Engineer in Artificial Intelligence from ESI-SBA, I am writing to express my strong enthusiasm for the PhD Fellowship in AI-Native Networking within the SYNAPSE project. My research background in developing physics-informed digital twins and optimized graph-based systems aligns perfectly with your vision of building synergetic network-AI platforms for mission-critical healthcare.

The core objective of SYNAPSE—integrating semantic hypergraph intelligence with distributed learning to recognize clinical urgency—is exactly the type of high-impact technical challenge I am prepared to tackle. My current thesis at Sonatrach focuses on resolving the "Latency-vs-Precision" trade-off for industrial digital twins. By implementing state-of-the-art **Mamba-SSM** and **xLSTM** architectures, I achieved a **0.04ms** inference time on edge devices, ensuring real-time responsiveness without compromising accuracy. This experience with temporal sequences and resource-constrained AI deployment is directly applicable to capturing network dynamics in AI-native ecosystems.

Furthermore, my work on **Graph Neural Networks (GNNs)** for structural anomaly detection (Career Connect and Elliptic++ projects) has equipped me with a deep understanding of relational intelligence. I am particularly excited about the prospect of developing **Semantic Hypergraph representations** to encode urgency propagation paths, as this mirrors the structural modeling I have performed for complex industrial networks.

Joining the Section for Medical ICT at OUS represents a unique opportunity to apply my expertise in GNNs and metaheuristic optimization (PSO/GA) to the critical domain of remote patient monitoring. I am a structured and adaptable researcher, committed to collaborative excellence and high-level scientific publication.

Thank you for your time and consideration. I look forward to the possibility of discussing how my technical background can contribute to the success of the SYNAPSE project.

Sincerely,

\bigskip

Ismail AISSAOUI

\end{document}
